\par\noindent\textbf{(1)} 归一入射线偏振 \(\boldsymbol e_{\rm in}=(\cos\alpha,\sin\alpha)^T\)。若 \(y\) 轴比 \(x\) 轴多延迟相位 \(\delta\)，波片 Jones矩阵可取 \(J(\delta)=\operatorname{diag}(1,e^{i\delta})\)，输出 \(\boldsymbol e_{\rm out}=(\cos\alpha,e^{i\delta}\sin\alpha)^T\)，整体公共相位无物理影响。
\par\noindent\textbf{(2)} 当 \(\alpha=\pi/4,\delta=\pi/2\) 时输出 \(\boldsymbol e=(1,i)^T/\sqrt2\)。两个正交分量幅相相等且相差四分之一周期，实场合成圆轨迹，满足 \(|E_x|=|E_y|\)，其旋向依时谐、观察方向与左右旋定义。
\par\noindent\textbf{(3)} 定义 \(S_0=|E_x|^2+|E_y|^2\)、\(S_1=|E_x|^2-|E_y|^2\)、\(S_2=2\operatorname{Re}(E_x^*E_y)\)、\(S_3=2\operatorname{Im}(E_x^*E_y)\)。输出有 \(S_0=1,S_1=\cos2\alpha,S_2=\sin2\alpha\cos\delta,S_3=\sin2\alpha\sin\delta\)，且 \(S_1^2+S_2^2+S_3^2=S_0^2\)。
\par\noindent\textbf{(4)} Jones矢量只能描述完全相干的确定偏振纯态。若光束在统计平均后出现混合偏振，须使用 \(2\times2\) 相干矩阵或Stokes–Mueller形式；此时一般满足 \(S_1^2+S_2^2+S_3^2<S_0^2\)，不能凭单一输出椭圆假定完全偏振。
